\chapter{Three Models for What Comes Next}
\label{ch:three-models}

\epigraph{The conversation is no longer about whether organisations restructure around AI, but how.}{}

\section{The Emerging Paradigms}

The consultancy world has converged on the fact that a structural shift is underway. McKinsey has explicitly adopted ``agentic organisation'' as a named paradigm, publishing ``Six Shifts to Build the Agentic Organization of the Future'' \parencite{McKinseySixShifts2026}. BCG published ``Unlocking the AI-First Organization: An Agentic Shift'' \parencite{BCGAgenticShift2025} and ``The Emerging Agentic Enterprise'' \parencite{BCGAgenticEnterprise2025}. The question is no longer whether organisations restructure around AI, but how --- and specifically, how much human coordination remains in the system.

Three distinct models have emerged. They are not competitors in the commercial sense. They describe different positions on a spectrum from full AI delegation to structured human-AI collaboration. Each carries different assumptions about the jagged frontier, the vigilance problem, and the irreducibility of human judgment.

\section{Model A: The Intelligence Layer}

\begin{figure}[h]
\centering
\fbox{\parbox{0.8\textwidth}{\centering\vspace{3cm}\textbf{[FIGURE: The Intelligence Layer Architecture]}\\\textit{Block's model: Company World Model, Customer World Model, and Intelligence Layer composing capabilities into solutions}\vspace{3cm}}}
\caption{Block's Intelligence Layer Model}
\label{fig:intelligence-layer}
\end{figure}

\subsection{Source and Core Thesis}

``From Hierarchy to Intelligence,'' published by Jack Dorsey and Roelof Botha on 31 March 2026 \parencite{DorseyBotha2026}, proposes the most radical of the three models: replace middle management entirely. Build a ``Company World Model'' --- an AI system continuously tracking what is being built, what is blocked, and where resources are allocated --- alongside a ``Customer World Model'' constructed from transaction data. An ``Intelligence Layer'' composes capabilities into solutions proactively. Three roles remain: Individual Contributors (deep specialists), Directly Responsible Individuals (who own cross-cutting problems for fixed periods), and Player-Coaches (who build things while developing people).

\subsection{The Case For}

The diagnosis is correct. The analysis of hierarchy as an information routing protocol is historically sound and analytically precise. The backing of Sequoia Capital lends credibility; their thesis that ``speed is the best predictor of startup success'' aligns with the argument that hierarchical latency is a competitive disadvantage. Block's stock surged 24 per cent on the announcement, and the company launched ``Managerbot'' --- a proactive AI agent described by VentureBeat as ``the clearest proof point yet.''

If an organisation possesses massive proprietary data (millions of transactions) and operates in a remote-first culture where work is already machine-readable, a world model is genuinely feasible as a coordination mechanism.

\subsection{The Case Against}

The paper was published one month after Block cut 4,000 of its 10,000-plus employees. A previous layoff memo from March 2025, eliminating 931 positions, \emph{explicitly stated} that the cuts were ``not about replacing people with AI.'' Eleven months later, AI was the entire thesis. \textcite{Zamost2026}, Block's former Head of Communications, observed that ``shrinking the policy team and eliminating diversity and inclusion roles looks like standard prioritisation, not an AI-driven reinvention.''

The operational evidence is sobering: 95 per cent of AI-generated code at Block still requires human modification, according to current and former employees. The comparison Block draws with Haier is structurally backwards: Haier \emph{decentralises} authority to micro-enterprises with profit-and-loss accountability; Block \emph{centralises} authority in an AI system and a single decision-maker.

\textcite{Spicer2026} observes that ``every wave of office technology produces delayering that subsequently reverses as complexity reasserts itself.'' The Klarna cautionary tale is instructive: the CEO claimed AI was performing the work of 700 customer service agents, but customer satisfaction dropped sharply. Klarna is now quietly rehiring humans, and the CEO admitted that ``cost unfortunately seems to have been a too predominant evaluation factor'' \parencite{Klarna2025}. An Orgvue/Forrester survey found that 55 per cent of companies that rushed to replace humans with AI now regret the decision \parencite{OrgvueForrester2025}.

\begin{cautionbox}
Block is right about the diagnosis and wrong about the prescription. The leap from ``information routing is automated'' to ``eliminate the humans who do coordination'' ignores the jagged frontier, the vigilance problem, the Klarna evidence, and the historical pattern of delayering reversal.
\end{cautionbox}

\section{Model B: The Agent Colony}

\begin{figure}[h]
\centering
\fbox{\parbox{0.8\textwidth}{\centering\vspace{3cm}\textbf{[FIGURE: The Agent Colony]}\\\textit{Every's model: personal agents developing specialisation, memory, and personality alongside their human partners}\vspace{3cm}}}
\caption{Every's Agent Colony Model}
\label{fig:agent-colony}
\end{figure}

\subsection{Source and Core Thesis}

The agent colony model emerges from the experiments at Every.to, documented by Dan Shipper in his ``Chain of Thought'' column and the podcast ``We Gave Every Employee an AI Agent'' \parencite{Shipper2026, Klaassen2025}. The core thesis: give every employee a persistent, personalised AI agent. The agents develop specialisation, memory, and personality that mirror their human partner. Organisational coordination emerges from cultural norms and public visibility rather than architectural redesign.

\subsection{What Works}

Several features of the Every experiment are genuinely instructive:

\textbf{Personal ownership creates accountability.} Each agent inherits its human's specialisation and reputation, creating a parallel organisational chart of specialised agents. Trust in an agent is derivative of trust in its human owner. As Shipper puts it: ``When R2C2 makes a mistake, I'm embarrassed --- because he's mine.'' The distinction between a personal agent and a generic tool matters: ``Claude is everybody's. A Plus One is mine.''

\textbf{Cultural norms do coordination work.} Public visibility in shared channels drives adoption and accountability. Agents share skills rapidly: ``They just write up a little document and send it. It's like in The Matrix when Neo says `I know kung fu.'\,''

\textbf{Specialisation emerges organically.} Growth questions go to one agent, product questions to another. The system self-organises around demonstrated capability rather than assigned responsibility.

\textbf{Management skill becomes universal.} If every employee has a personal agent, they must onboard, delegate to, evaluate, and coach it. The skill of management does not disappear; it democratises. This directly contradicts Block's thesis that management should be eliminated.

\subsection{What Does Not Work}

\textbf{The ant death spiral.} If one agent messages a channel monitored by other agents and the configurations are imprecise, the agents will respond to each other indefinitely --- ``burning millions of tokens'' until a human intervenes. This is a failure of coordination at the system level, precisely the function that hierarchy was designed to provide.

\textbf{Error multiplication.} Naive multi-agent systems exhibit 17-times error multiplication \parencite{TDS2026}, where errors in one agent's output propagate and amplify through downstream agents.

\textbf{Scale limitations.} Every's experiment involved approximately twenty people. Cultural self-organisation works at this scale. At two hundred people, you need structure. At two thousand, you need architecture. The experiment is an existence proof of a dynamic that has not yet been shown to generalise.

\section{Model C: The Dynamic Agentic Mesh}

\begin{figure}[h]
\centering
\fbox{\parbox{0.8\textwidth}{\centering\vspace{3cm}\textbf{[FIGURE: The Dynamic Agentic Mesh Architecture]}\\\textit{DreamLab's model: three-layer architecture with Discovery Engine, Orchestration Layer, and Declarative Governance}\vspace{3cm}}}
\caption{DreamLab's Dynamic Agentic Mesh}
\label{fig:agentic-mesh}
\end{figure}

\subsection{Core Thesis}

The Dynamic Agentic Mesh retains human judgment at intersection points. Middle managers evolve from information routers to \meshterm{judgment brokers} --- federated guardians of a mesh where agentic orchestration handles information flow and humans handle what AI cannot: navigating the jagged frontier, maintaining vigilance, ethical adjudication, and cultural coherence.

\subsection{Three Layers}

The mesh operates through three integrated layers.

\textbf{The Discovery Engine} captures bottom-up insight. Passive agent monitoring of communication channels surfaces shadow workflows --- the unofficial optimisations that staff have already invented. Natural-language Insight Portals allow anyone to propose cross-functional shortcuts. Tacit knowledge is automatically codified into agentic instructions.

\textbf{The Orchestration Layer} provides agentic coordination. Directed Acyclic Graphs re-route tasks in real time based on human and agent availability. Data transforms from ``dashboard view'' to ``data product,'' carrying its own service-level agreements, readiness checks, and freshness guarantees. Specialist agents execute defined skill domains with ontological precision.

\textbf{Declarative Governance} embeds governance as code. AI-enforced policies replace manual stewardship. Bias detection, security protocols, and value-alignment checks are woven into every task transition. The judgment broker reviews only genuine edge cases and strategic decisions that exceed agent authority.

\subsection{The Insight Ingestion Loop}

The mesh converts individual discoveries into organisational capability through a five-stage loop:

\begin{enumerate}
  \item \textbf{Discovery}: An employee finds a shortcut. The mesh detects the pattern.
  \item \textbf{Codification}: The system maps the new path as a proposed DAG, formalised with provenance and dependency checks.
  \item \textbf{Validation}: The judgment broker reviews for strategic fit, bias risk, and compliance.
  \item \textbf{Integration}: The validated workflow is promoted into the live mesh with SLAs and ownership.
  \item \textbf{Amplification}: The mesh propagates the pattern to other teams. Mesh velocity --- insight-to-integration time --- becomes a strategic KPI.
\end{enumerate}

\subsection{Honest Limitations}

The DreamLab Mesh has been demonstrated at 15-person scale in a creative technology context, led by a founder who is also the system architect. The honest questions are:

\begin{enumerate}
  \item Does the judgment broker role work when the broker did not build the mesh?
  \item Does the insight ingestion loop work in regulated industries where shadow workflows are compliance violations?
  \item Does the OWL~2 ontology scale to organisations with thousands of concepts?
  \item Can middle managers be trained in judgment brokerage fast enough to matter?
\end{enumerate}

These are genuine open questions. Presenting them honestly is more credible than concealing them.

\section{Comparative Analysis}

\begin{table}[h]
\centering
\caption{Comparative analysis of three post-hierarchical models}
\label{tab:three-models}
\begin{tabularx}{\textwidth}{l L{3cm} L{3cm} X}
\toprule
\textbf{Dimension} & \textbf{Block} & \textbf{Every} & \textbf{DreamLab Mesh} \\
\midrule
Coordination & AI world model & Cultural norms & Declarative governance + DAG orchestration \\
\addlinespace
Management & Eliminate & Everyone manages their agent & Judgment brokers at intersections \\
\addlinespace
Scale (proven) & 10,000+ (intended) & $\sim$20 & 15, designed for 50--500 \\
\addlinespace
Evidence & 95\% code needs humans; Klarna reversed & Works at small scale; ant death spiral & VisionClaw running; OWL~2 provenance \\
\addlinespace
Collaboration mode & Self-Automator org & Cyborg org & Centaur org \\
\addlinespace
Primary risk & Falling asleep at the wheel & 17x error multiplication at scale & Unproven beyond 15-person team \\
\bottomrule
\end{tabularx}
\end{table}

The evidence from Chapter~\ref{ch:jagged-frontier} indicates that Self-Automator organisations --- those that delegate wholesale to AI --- produce the weakest outcomes. This is the empirical case against Block's model and for organisational designs that maintain active human engagement at the frontier boundary.
